\documentclass[11pt,letterpaper]{article}
\usepackage[margin=1in]{geometry}
\usepackage[T1]{fontenc}
\usepackage{lmodern}
\usepackage{microtype}
\usepackage{amsmath,amssymb}
\usepackage{booktabs,array,enumitem,xcolor}
\usepackage{titlesec}
\usepackage[colorlinks=true,linkcolor=black,urlcolor=blue!50!black,citecolor=black]{hyperref}
\setlist{itemsep=2pt,topsep=3pt}
\setlength{\parindent}{0pt}\setlength{\parskip}{6pt}
\titleformat{\section}{\normalfont\large\bfseries}{\thesection}{0.6em}{}
\titlespacing*{\section}{0pt}{12pt}{4pt}
\pagestyle{plain}
\newcommand{\plainwords}[1]{\par\medskip\noindent\fcolorbox{black!25}{black!4}{\parbox{\dimexpr\linewidth-2\fboxsep-2\fboxrule}{\textbf{In plain words.} #1}}\par\medskip}
\newcommand{\titleblock}[3]{%
  {\LARGE\bfseries #1\par}\vspace{4pt}
  {\large #2\par}\vspace{2pt}
  {\small\color{black!70} Shakil A.\ Rafi, Ph.D.\ \textbar\ University of Arkansas \textbar\ #3\par}\vspace{2pt}
  \hrule\vspace{6pt}}

\begin{document}
\titleblock{RustQC}{A fast, memory-bounded FASTQ quality-control tool in Rust}{Software write-up, November 2026}
\textit{Source: \href{https://github.com/2shakilrafi/rustqc}{github.com/2shakilrafi/rustqc} (MIT). Status: v0.4.0; conda-forge publication in review.}

\plainwords{Before anyone analyzes sequencing data, they check its quality, and FastQC is the standard tool for that. I rebuilt its analyses in Rust so the checks run in one pass with a fixed memory ceiling.}

\section{Introduction}
Quality control is the first step of most sequencing pipelines. RustQC targets FastQC 0.12.1: it implements the eleven active analysis modules and writes self-contained reports plus FastQC-compatible data files, so existing pipeline tooling can read its output. It is command-line first and does not reproduce FastQC's desktop interface.

\section{Design}
\begin{itemize}
\item \textbf{One streaming pass} over each input, with only a small bounded number of record batches in flight.
\item \textbf{Bounded state.} Fixed quality histograms replace per-base quality vectors, duplicate tracking stops at 100{,}000 unique sequence prefixes, and per-tile analysis follows FastQC's sampling rule. Memory therefore grows with the longest read, not the number of reads.
\item \textbf{Parallelism.} A single file uses a parser, workers, and an ordered reduction; multiple files run concurrently without oversubscribing the machine.
\item \textbf{Compatibility.} Default thresholds, position grouping, and GC interpolation follow FastQC; outputs include an interactive HTML report, a FastQC-layout zip, and structured JSON.
\end{itemize}

\section{Validation and results}
Correctness is checked with a Rust test suite and formatting and lint gates (\texttt{cargo fmt}, \texttt{cargo test}, \texttt{cargo clippy} with warnings denied). My early in-memory benchmarks on FASTX input were more than $40\times$ faster than FastQC. I treat that as a first signal only: in-memory timings leave out compressed-file decompression, and I have not yet published an end-to-end comparison that includes gzip input, peak memory, and thread scaling. Writing that benchmark, with a head-to-head output comparison against FastQC on shared data, is the next step.

\section{Discussion}
The aim was parity with FastQC's results, not a new method, so the engineering question is whether the outputs agree and the speed holds up on realistic files. For distribution, the package is installable with Cargo today. I have submitted a conda-forge recipe under the name \texttt{rustqc-fastq}, because an unrelated project already uses \texttt{rustqc} on Bioconda. That review is still open.

\textbf{Contact.} sarafi@uark.edu
\end{document}
